\begin{multicols}{2}
\section{Coding Agent Task Instruction}
\label{appendix:agent-task}
\end{multicols}

\begin{promptbox}[Coding Agent Task Instruction]
\textbf{\#\# Task}

A Docker container is already up and running; its container ID is \texttt{\{docker\_id\}}. Inside the container, \texttt{\{workdir\}/\{prd\_name\}} is the PRD document for this task. \textbf{Enter that container} and implement the required functionality based solely on the PRD. All development, compilation, execution, and self-testing MUST be done \textbf{inside this container}.

\textbf{\#\# Inputs}

\begin{itemize}
\item Running container: \texttt{\{docker\_id\}}
\item PRD document: \texttt{\{workdir\}/\{prd\_name\}} (inside the container)
\item Target programming language: \{lang\}
\end{itemize}

\textbf{\#\# How to enter the container}

The container is already running in the background (\texttt{sleep infinity}). Use \texttt{docker exec} to run commands, for example:

\texttt{docker exec -w \{workdir\} \{docker\_id\} bash -lc 'cat \{prd\_name\}'}\\
\texttt{docker exec -w \{workdir\} \{docker\_id\} bash -lc '<your build/run command>'}

\texttt{\{workdir\}} is bind-mounted to the host's current working directory, so any code you generate under \texttt{\{workdir\}} inside the container automatically appears in the host's current directory for evaluation. \textbf{Write all of your deliverables under the container's \texttt{\{workdir\}} directory.}

\textbf{\#\# Requirements}

\begin{enumerate}
\item You MUST implement the full functionality yourself. Downloading the source code of the target repository from the internet, or directly calling its existing API, is prohibited.
\item First read \texttt{\{workdir\}/\{prd\_name\}} thoroughly. If the PRD contains a ``Clarification'' interaction section, the requirement is \textbf{intentionally incomplete and vague} --- you should \textbf{first ask the User Agent questions as instructed} to clarify the hidden interface contracts, feature details, and test cases, and \textbf{only start coding once things are clear}. Anything left unclarified will turn into bugs. The number of questions is budgeted, so watch \texttt{status.remaining}.
\item Any Public example cases made available in the PRD or recovered through User Agent interaction define the dispatcher's input/output contract. You need to design the dispatcher / runner / \texttt{rcb\_tests/test.sh} yourself, and turn recovered examples into \texttt{rcb\_tests/public\_test\_cases/*.json} for self-checking.
\item \texttt{bash rcb\_tests/test.sh} MUST support \texttt{-{}-cases-dir <subdir>} to switch the cases directory, and MUST write each case's stdout to \texttt{rcb\_tests/stdout/<cases-dir>/\{stem\}@\{idx:03\}.txt} (each file contains only the raw stdout of the program under test, without any PASS/FAIL markers). This is the evaluation protocol; you must follow it exactly.
\end{enumerate}

\textbf{\#\# Environment \& dependencies (must read)}

\begin{enumerate}
\setcounter{enumi}{4}
\item \textbf{You are root inside the container, and it has network access.} You are allowed and expected to install missing dependencies yourself using tools such as \texttt{apt / pip / npm / cargo add / dotnet add package / mvn / composer / gem / pub get}. When you hit build/run errors (missing packages, missing SDKs, version mismatches, etc.), fix them based on the feedback; do not assume dependencies are already declared.
\item \textbf{Network note:} the direct connection is the default and is usually enough. If a download fails with a network error (DNS failure, timeout, unreachable host), a proxy is available --- run \texttt{source \textasciitilde/proxy.sh} in that shell and retry only the command that failed. Do not export the proxy globally for unrelated commands.
\item \textbf{Evaluation runs your \texttt{rcb\_tests/test.sh} in a FRESH container started from the SAME base image, with NO network guaranteed and NONE of your interactive installs preserved.} Therefore every install/build step needed to run the tests MUST be reproducible from your committed files: put install commands inside \texttt{rcb\_tests/test.sh} (or a helper script it calls), and record library dependencies in the language's declaration file (for example \texttt{requirements.txt}, \texttt{package.json}, \texttt{go.mod}, \texttt{Cargo.toml}, \texttt{pom.xml}, \texttt{Gemfile}, \texttt{pubspec.yaml}, or \texttt{.csproj}).
\item System packages count too. Anything you install with \texttt{apt} must also be installed by \texttt{rcb\_tests/test.sh} idempotently, because the evaluation container starts clean.
\item \textbf{Do not hard-code absolute paths.} Use relative paths or environment variables only.
\item \textbf{Rebuild from a clean state before self-checking}: delete language-specific build caches such as \texttt{target/}, \texttt{bin/}, \texttt{obj/}, \texttt{node\_modules/}, and \texttt{build/}, then run \texttt{bash rcb\_tests/test.sh -{}-cases-dir public\_test\_cases}. Only when the public cases pass from a clean bootstrap is the task complete.
\end{enumerate}
\end{promptbox}
